% GENERATED FILE. Edit canonical lessons, then run npm run generate:books.
% canonical-source-hash: fnv1a64:c21aea19af74510c
% canonical-lessons: JA-C233-asobu, JA-C233-oyogu, JA-C233-suiei, JA-C233-hashiru, JA-C233-undou

\chapter{Swimming and Running}
\label{ch:ja-swimming-and-running}
\hlchaptermodality{233}

\begin{chapteropening}
\textbf{By the end of this chapter:} \emph{I can say how I played and exercised.}

Complete the last lesson of chapter 233: I can say how I played and exercised.
\end{chapteropening}

\section[\texorpdfstring{\ja{あそぶ} (asobu)}{あそぶ (asobu)}]{\ja{あそぶ} (asobu) — to play; to have fun}

\label{lesson:JA-C233-asobu}



\begin{quote}
Before the new one: say the Japanese for to refuse, then the Japanese for to inform.
\end{quote}

\subsection*{You'll want to know: \ja{あそぶ}}
\textbf{\ja{あそぶ}} — \emph{asobu} — \textquotedblleft{}to play; to have fun\textquotedblright{}.

\begin{quote}
\textbf{\ja{きのう} \ja{こうえんで} \ja{ともだちと} \ja{あそびました。}} — \emph{kinō kōen de tomodachi to asobimashita} — \textquotedblleft{}Yesterday I played with friends in the park.\textquotedblright{}
\end{quote}

Polite \textbf{\ja{あそびます}}, \emph{asobimasu}; a godan verb.

\begin{cousinweb}
\textbf{\ja{あそぶ}} (\emph{asobu}) is a native Japanese word (wago), inherited rather than borrowed from Chinese; its kanji was attached later as a kun'yomi (JMdict, KANJIDIC2).
\end{cousinweb}

\begin{grammarlens}[title={What you've built}]
One more word for everyday life at A2.
\end{grammarlens}

\subsection*{Your turn}
\begin{itemize}
\raggedright
  \item \emph{Say it:} \emph{asobu}
  \item \emph{Say it:} \emph{asobu}, once more
  \item \emph{Say it:} \emph{kinō kōen de tomodachi to asobimashita} — Yesterday I played with friends in the park.
  \item \emph{Recall:} say the Japanese for to refuse, then the Japanese for to inform, then say \emph{asobu} again
\end{itemize}

\begin{tcolorbox}[breakable,colback=teal!4,colframe=teal!35!black,title={Before you move on}]
What does \ja{あそぶ} mean? (\textquotedblleft{}To play; to have fun\textquotedblright{}.) Say it once more.
\end{tcolorbox}
\section[\texorpdfstring{\ja{およぐ} (oyogu)}{およぐ (oyogu)}]{\ja{およぐ} (oyogu) — to swim}

\label{lesson:JA-C233-oyogu}



\begin{quote}
Before the new one: say the Japanese for to inform, then the Japanese for to play.
\end{quote}

\subsection*{You'll want to know: \ja{およぐ}}
\textbf{\ja{およぐ}} — \emph{oyogu} — \textquotedblleft{}to swim\textquotedblright{}.

\begin{quote}
\textbf{\ja{せんしゅう} \ja{うみで} \ja{およぎました。}} — \emph{senshū umi de oyogimashita} — \textquotedblleft{}Last week I swam in the sea.\textquotedblright{}
\end{quote}

Polite \textbf{\ja{およぎます}}, \emph{oyogimasu}; a godan verb.

\begin{cousinweb}
\textbf{\ja{およぐ}} (\emph{oyogu}) is a native Japanese word (wago), inherited rather than borrowed from Chinese; its kanji was attached later as a kun'yomi (JMdict, KANJIDIC2).
\end{cousinweb}

\begin{grammarlens}[title={What you've built}]
One more word for everyday life at A2.
\end{grammarlens}

\subsection*{Your turn}
\begin{itemize}
\raggedright
  \item \emph{Say it:} \emph{oyogu}
  \item \emph{Say it:} \emph{oyogu}, once more
  \item \emph{Say it:} \emph{senshū umi de oyogimashita} — Last week I swam in the sea.
  \item \emph{Recall:} say the Japanese for to inform, then the Japanese for to play, then say \emph{oyogu} again
\end{itemize}

\begin{tcolorbox}[breakable,colback=teal!4,colframe=teal!35!black,title={Before you move on}]
What does \ja{およぐ} mean? (\textquotedblleft{}To swim\textquotedblright{}.) Say it once more.
\end{tcolorbox}
\section[\texorpdfstring{\ja{すいえい} (suiei)}{すいえい (suiei)}]{\ja{すいえい} (suiei) — swimming; swimming as a sport}

\label{lesson:JA-C233-suiei}



\begin{quote}
Before the new one: say the Japanese for to play, then the Japanese for to swim.
\end{quote}

\subsection*{You'll want to know: \ja{すいえい}}
\textbf{\ja{すいえい}} — \emph{suiei} — \textquotedblleft{}swimming; swimming as a sport\textquotedblright{}.

\begin{quote}
\textbf{\ja{なつやすみに} \ja{すいえいを} \ja{ならいました。}} — \emph{natsuyasumi ni suiei o naraimashita} — \textquotedblleft{}I took swimming lessons in the summer holidays.\textquotedblright{}
\end{quote}

\begin{cousinweb}
\textbf{\ja{すいえい}} (\emph{suiei}) is Sino-Japanese (kango): two Chinese characters read with their on'yomi, readings borrowed from Chinese. They are \emph{sui}, \textquotedblleft{}water\textquotedblright{} (Middle Chinese \emph{sywijX}) and \emph{ei}, \textquotedblleft{}swim\textquotedblright{} (\emph{hjwaengH}) (KANJIDIC2; Middle Chinese in Baxter's transcription).
\end{cousinweb}

\begin{grammarlens}[title={What you've built}]
One more word for everyday life at A2.
\end{grammarlens}

\subsection*{Your turn}
\begin{itemize}
\raggedright
  \item \emph{Say it:} \emph{suiei}
  \item \emph{Say it:} \emph{suiei}, once more
  \item \emph{Say it:} \emph{natsuyasumi ni suiei o naraimashita} — I took swimming lessons in the summer holidays.
  \item \emph{Recall:} say the Japanese for to play, then the Japanese for to swim, then say \emph{suiei} again
\end{itemize}

\begin{tcolorbox}[breakable,colback=teal!4,colframe=teal!35!black,title={Before you move on}]
What does \ja{すいえい} mean? (\textquotedblleft{}Swimming; swimming as a sport\textquotedblright{}.) Say it once more.
\end{tcolorbox}
\section[\texorpdfstring{\ja{はしる} (hashiru)}{はしる (hashiru)}]{\ja{はしる} (hashiru) — to run}

\label{lesson:JA-C233-hashiru}



\begin{quote}
Before the new one: say the Japanese for to swim, then the Japanese for swimming.
\end{quote}

\subsection*{You'll want to know: \ja{はしる}}
\textbf{\ja{はしる}} — \emph{hashiru} — \textquotedblleft{}to run\textquotedblright{}.

\begin{quote}
\textbf{\ja{けさ} \ja{こうえんを} \ja{はしりました。}} — \emph{kesa kōen o hashirimashita} — \textquotedblleft{}This morning I ran in the park.\textquotedblright{}
\end{quote}

Polite \textbf{\ja{はしります}}, \emph{hashirimasu}; a godan verb.

\begin{cousinweb}
\textbf{\ja{はしる}} (\emph{hashiru}) is a native Japanese word (wago), inherited rather than borrowed from Chinese; its kanji was attached later as a kun'yomi (JMdict, KANJIDIC2).
\end{cousinweb}

\begin{grammarlens}[title={What you've built}]
One more word for everyday life at A2.
\end{grammarlens}

\subsection*{Your turn}
\begin{itemize}
\raggedright
  \item \emph{Say it:} \emph{hashiru}
  \item \emph{Say it:} \emph{hashiru}, once more
  \item \emph{Say it:} \emph{kesa kōen o hashirimashita} — This morning I ran in the park.
  \item \emph{Recall:} say the Japanese for to swim, then the Japanese for swimming, then say \emph{hashiru} again
\end{itemize}

\begin{tcolorbox}[breakable,colback=teal!4,colframe=teal!35!black,title={Before you move on}]
What does \ja{はしる} mean? (\textquotedblleft{}To run\textquotedblright{}.) Say it once more.
\end{tcolorbox}
\section[\texorpdfstring{\ja{うんどう} (undō)}{うんどう (undō)}]{\ja{うんどう} (undō) — exercise; physical training}

\label{lesson:JA-C233-undou}



\begin{quote}
Before the new one: say the Japanese for swimming, then the Japanese for to run.
\end{quote}

\subsection*{You'll want to know: \ja{うんどう}}
\textbf{\ja{うんどう}} — \emph{undō} — \textquotedblleft{}exercise; physical training\textquotedblright{}.

\begin{quote}
\textbf{\ja{きのう} \ja{ちちと} \ja{うんどうを} \ja{しました。}} — \emph{kinō chichi to undō o shimashita} — \textquotedblleft{}Yesterday I exercised with my father.\textquotedblright{}
\end{quote}

\begin{cousinweb}
\textbf{\ja{うんどう}} (\emph{undō}) is Sino-Japanese (kango): two Chinese characters read with their on'yomi, readings borrowed from Chinese. They are \emph{un}, \textquotedblleft{}carry, transport\textquotedblright{} (Middle Chinese \emph{hjunH}) and \emph{dō}, \textquotedblleft{}move\textquotedblright{} (\emph{duwngX}) (KANJIDIC2; Middle Chinese in Baxter's transcription).
\end{cousinweb}

\begin{grammarlens}[title={What you've built}]
One more word for everyday life at A2.
\end{grammarlens}

\subsection*{Your turn}
\begin{itemize}
\raggedright
  \item \emph{Say it:} \emph{undō}
  \item \emph{Say it:} \emph{undō}, once more
  \item \emph{Say it:} \emph{kinō chichi to undō o shimashita} — Yesterday I exercised with my father.
  \item \emph{Recall:} say the Japanese for swimming, then the Japanese for to run, then say \emph{undō} again
\end{itemize}

\begin{tcolorbox}[breakable,colback=teal!4,colframe=teal!35!black,title={Before you move on}]
What does \ja{うんどう} mean? (\textquotedblleft{}Exercise; physical training\textquotedblright{}.) Say it once more.
\end{tcolorbox}
